% A finding kept out of the paper (user's decision, 30 Oct 2026): FedECS-LT with its trained generator replaced by class
% Gaussians from exact statistics. Every number is a slot of docs/papers/generated/numbers.tex (note/... and abl/...),
% written by docs/papers/fill.py. Build: pdflatex gaussian_sampler.tex (two passes).
\documentclass[11pt,a4paper]{article}
\usepackage[T1]{fontenc}
\usepackage{lmodern}
\usepackage[a4paper,margin=2.7cm]{geometry}
\usepackage{microtype}
\usepackage{amsmath,amssymb}
\usepackage{graphicx}
\usepackage[export]{adjustbox}
\usepackage{booktabs}
\usepackage{xcolor,etoolbox,xspace}
\usepackage{caption}
\captionsetup{font=small,labelfont=bf}
\usepackage[colorlinks,linkcolor=black,citecolor=black,urlcolor=black]{hyperref}
\newcommand{\genpath}{../papers/generated}
% Shared macros of both versions of the paper (elsevier/main.tex, ieee/main.tex).
%
% RESULT SLOTS. Every number, table and figure that comes from an experiment is a slot. A slot shows a red
% placeholder until docs/papers/fill.py has written its value into docs/papers/generated/; then the value appears by
% itself. Never type a result into the text: run fill.py and rebuild (docs/papers/GUIDE.md).
%   \R{key}                    a number (generated/numbers.tex), e.g. \R{main/cifar10/FedECS-LT/bal}
%   \slottab{name}{fallback}   a tabular from generated/tables/name.tex; `fallback' (a placeholder tabular) until then
%   \slotfig{name}{size}{what} a figure from generated/figures/name.pdf; size = col (one column) or wide (page)
%   \claim{text}               a sentence whose wording depends on results: blue in draft mode, check it (GUIDE.md)
% \finalslotstrue in a wrapper turns every remaining placeholder into a LaTeX error, so that nothing unfilled is
% submitted.
\newif\iffinalslots
\makeatletter
\emergencystretch=2em   % placeholder names cannot be hyphenated
\providecommand{\genpath}{../generated}

\newcommand{\slotmissing}[1]{%
  \iffinalslots\PackageError{slots}{Result slot `#1' is not filled}{Run docs/papers/fill.py.}\fi
  \textcolor{red}{\textbf{[\detokenize{#1}]}}}

% numbers: generated/numbers.tex holds lines \defnum{key}{value}
\newcommand{\defnum}[2]{\expandafter\def\csname num@#1\endcsname{#2}}
\newcommand{\R}[1]{\ifcsname num@#1\endcsname\csname num@#1\endcsname\else\slotmissing{#1}\fi}
\InputIfFileExists{\genpath/numbers.tex}{}{}

% tables: the generated tabular, scaled down only if it is wider than the line
\newcommand{\slottab}[2]{%
  \IfFileExists{\genpath/tables/#1.tex}%
    {\begin{adjustbox}{max width=\linewidth}\input{\genpath/tables/#1.tex}\end{adjustbox}}%
    {\iffinalslots\PackageError{slots}{Table slot `#1' is not filled}{Run docs/papers/fill.py.}\fi
     \begin{adjustbox}{max width=\linewidth}#2\end{adjustbox}\par
     {\scriptsize\textcolor{red}{Placeholder: table slot \texttt{\detokenize{#1}} (GUIDE.md)}}}}
% one placeholder row: method name and n empty cells (the cells are built first, then emitted, so that the &s act
% as column separators of the table)
\newcommand{\phrow}[2]{\begingroup\count@=#2\relax\def\ph@c{}%
  \loop\ifnum\count@>\z@\edef\ph@c{\unexpanded\expandafter{\ph@c & --}}\advance\count@\m@ne\repeat
  \expandafter\endgroup\expandafter\ph@emit\expandafter{\ph@c}{#1}}
\newcommand{\ph@emit}[2]{#2#1\\}

% figures: the drawn PDF at its print size (column 3.5 in, page 7.16 in; IEEE format), never wider than the line
\newlength{\slotfigh}
\newcommand{\slotfig}[3]{%
  \IfFileExists{\genpath/figures/#1.pdf}%
    {\includegraphics[max width=\linewidth]{\genpath/figures/#1.pdf}}%
    {\iffinalslots\PackageError{slots}{Figure slot `#1' is not filled}{Run docs/papers/fill.py --figures.}\fi
     \ifstrequal{#2}{wide}{\setlength{\slotfigh}{0.28\textwidth}}{\setlength{\slotfigh}{0.55\linewidth}}%
     \fbox{\parbox[c][\slotfigh][c]{\dimexpr\linewidth-2\fboxsep-2\fboxrule\relax}{\centering\small
        \textcolor{red}{\textbf{FIGURE PLACEHOLDER: \texttt{\detokenize{#1}}}}\\[2pt]\footnotesize #3\par}}}}

% parts that exist only when their experiment came back: in draft mode always shown (with placeholders), in the
% final build (\finalslotstrue) left out when the result is missing, so that the paper never shows an empty slot.
%   \whentab{name}{...}  \whenfig{name}{...}  \whennum{key}{...}  \whenanytab{a,b,c}{...}
\newcommand{\whentab}[2]{\iffinalslots\IfFileExists{\genpath/tables/#1.tex}{#2}{}\else#2\fi}
\newcommand{\whenfig}[2]{\iffinalslots\IfFileExists{\genpath/figures/#1.pdf}{#2}{}\else#2\fi}
\newcommand{\whennum}[2]{\iffinalslots\ifcsname num@#1\endcsname#2\fi\else#2\fi}
\newif\ifslot@any
\newcommand{\whenanytab}[2]{\iffinalslots\slot@anyfalse
  \@for\slot@x:=#1\do{\IfFileExists{\genpath/tables/\slot@x.tex}{\global\slot@anytrue}{}}%
  \ifslot@any#2\fi\else#2\fi}

% a sentence to check against the results before submission (blue in draft mode, plain in final mode)
\newcommand{\claim}[1]{\iffinalslots#1\else\textcolor{blue}{#1}\fi}

% notation
\newcommand{\E}{\mathbb{E}}
\newcommand{\Rset}{\mathbb{R}}
\newcommand{\LSE}{\mathrm{LSE}}
\newcommand{\KL}{\mathrm{KL}}
\newcommand{\ER}{ER\xspace}
\newcommand{\ours}{FedECS-LT\xspace}
\newcommand{\geflf}{GeFL-F\xspace}
\newcommand{\pos}[1]{\textbf{#1}}
\makeatother

\finalslotstrue
\setlength{\parskip}{0.35em}
\linespread{1.08}
\renewcommand{\paragraph}[1]{\par\medskip\noindent\textbf{#1.}\enspace}
\newcommand{\G}[1]{\R{note/cifar10/w/ Gaussian sampler/#1}}
\newcommand{\Gc}[1]{\R{note/cifar100/w/ Gaussian sampler/#1}}

\begin{document}
\begin{center}
{\Large\bfseries FedECS-LT without a trained generator}\\[4pt]
{\small A finding from the component ablation, kept out of the paper}
\end{center}
\medskip

\paragraph{The question} The component ablation of FedECS-LT includes one run that the paper does not report: the
trained generator of Stage II replaced by class Gaussians built from exact statistics. Everything else is as in
FedECS-LT (consensus soft labels, logit adjustment, ten synthetic epochs, the exact readout). Reviewer E asked for this
arm: if it matches FedECS-LT, Stage II is unnecessary. It does better than matching.

\paragraph{What the sampler is} After Stage I the clients send, besides the class counts and feature sums, the sum of
$hh^\top$ over their features. From these the server forms the class means $\mu_c$ and one within-class covariance
$\Sigma$ shared by all classes, shrunk toward a multiple of the identity with $\gamma=0.1$, the same estimator as the
readout's (\texttt{fedecs.gaussian\_from\_sums}). A class-$c$ feature is then $\mathrm{ReLU}(\mu_c+L\varepsilon)$ with
$\varepsilon\sim\mathcal N(0,I)$ and $LL^\top=\Sigma$ (\texttt{models.GaussFeatureGen}). There is no Stage II, no class
row, no gate and no moment correction: the sampler has the exact mean and, before the ReLU, the exact pooled
covariance by construction. It is the sampler of CCVR, applied to the generator-based pipeline.

\paragraph{The results} Table~\ref{tab:g} compares it with the full FedECS-LT on the same splits and seeds (CIFAR-10-LT
and CIFAR-100-LT, imbalance ratio 100, ten clients, one per architecture).

\begin{table}[h]
\centering\small
\caption{FedECS-LT with its trained generator and with the Gaussian sampler: balanced accuracy (\%) of the heads as
trained and with the exact readout, mean $\pm$ standard deviation over the seeds that finished; the differences are
the sampler minus the full method, paired over common seeds, with uncorrected $p$; Holm: corrected over the seven
component changes of the ablation.}
\label{tab:g}
\begin{tabular}{lcc}\toprule
& CIFAR-10-LT & CIFAR-100-LT\\\midrule
Seeds & \G{seeds} & \Gc{seeds}\\
FedECS-LT, heads as trained & \R{main/cifar10/FedECS-LT (heads only)/bal} & \R{main/cifar100/FedECS-LT (heads only)/bal}\\
Gaussian sampler, heads as trained & \G{final_bal} & \Gc{final_bal}\\
\quad difference ($p$) & \G{diff} (\G{p}) & \Gc{diff} (\Gc{p})\\
FedECS-LT with the readout & \R{main/cifar10/FedECS-LT/bal} & \R{main/cifar100/FedECS-LT/bal}\\
Gaussian sampler with the readout & \G{er_bal} & \Gc{er_bal}\\
\quad difference ($p$; Holm) & \G{er_diff} (\G{er_p}; \G{er_p_holm}) & \Gc{er_diff} (\Gc{er_p}; \Gc{er_p_holm})\\
\quad tail recall difference ($p$) & \G{er_tail_diff} (\G{er_tail_p}) & \Gc{er_tail_diff} (\Gc{er_tail_p})\\\midrule
Communication per client (MB) & \G{comm_mb_per_client} vs \G{comm_mb_per_client_ours} &
\Gc{comm_mb_per_client} vs \Gc{comm_mb_per_client_ours}\\
Wall-clock time of a run (s) & \G{time_s} vs \G{time_s_ours} & \Gc{time_s} vs \Gc{time_s_ours}\\\bottomrule
\end{tabular}
\end{table}

On CIFAR-10-LT the sampler is \G{er_diff} points ahead with the readout, significant after Holm's correction over the
seven component changes ($p=\G{er_p_holm}$) and above the noise floor of \R{repeat/er/floor} points, and \G{diff}
points ahead with heads as trained. Its tail recall with the readout is \G{er_tail_diff} points higher. On
CIFAR-100-LT it is \Gc{er_diff} points ahead, from \Gc{seeds} seeds, which is unresolved. It needs
\G{comm_mb_per_client_ratio} times less communication and \G{time_s_ratio} times less time on CIFAR-10-LT
(\Gc{comm_mb_per_client_ratio} and \Gc{time_s_ratio} on CIFAR-100-LT), because Stage II, two hundred rounds of FedAvg
on the generator, is the largest part of FedECS-LT's cost.

The same holds for GeFL-F. The paper's variants table already contains GeFL-F with its generator replaced by the
same Gaussians, with logit adjustment. That change raises GeFL-F by
\R{abl/cifar10/GeFL-F w/ Gaussian generator + LA/paired} points as trained on CIFAR-10-LT
($p=\R{abl/cifar10/GeFL-F w/ Gaussian generator + LA/p}$), its tail recall by
\R{abl/cifar10/GeFL-F w/ Gaussian generator + LA/tail_paired}, and its accuracy with the readout by
\R{abl/cifar10/GeFL-F w/ Gaussian generator + LA/er_paired} ($p=\R{abl/cifar10/GeFL-F w/ Gaussian generator + LA/er_p}$).

\paragraph{Why it may work} The paper's third diagnosis is that GeFL-F's generated features do not carry their
label: the referee recognises them near chance (\R{main/cifar10/GeFL-F/fidelity_all}\,\% on CIFAR-10-LT, where chance is 10\%). A Gaussian
centred on the exact class mean carries its label by construction, as far as the means differ, and its spread is the
real one. The heads then learn the classes a client lacks from features that look like them on average, which is what
the generator was meant to provide. That a two-moment sampler beats a trained conditional VAE suggests that, on this
shallow frozen feature space, the generator adds nothing that the first two moments do not already carry, and that
its training loses part of what they carry. This is an interpretation; the runs do not test it.

\paragraph{What it would mean for the paper} If the sampler became the method, FedECS-LT would be simpler and
cheaper: exact class statistics, a Gaussian sampler, consensus soft labels, logit adjustment and the exact readout.
The diagnosis would stay as it is; the repair of the generator (anchor, gate, CCWA, moment correction) would be
replaced by not training a generator at all, which fits the method's name better. The anchored generator would
become an ablation.

\paragraph{What it would cost to adopt} New FedECS-LT runs with the sampler on every dataset where the method is
reported: MNIST, Fashion-MNIST and SVHN; the six naturally imbalanced datasets; CIFAR-100-LT seed 5; the new splits
(seeds 6--8); and the 50-client setting. Without Stage II a CIFAR run took \G{time_s} s (CIFAR-10-LT) and \Gc{time_s} s (CIFAR-100-LT) on a shared Kaggle GPU. The baselines need no new
runs. The method section, the overview figure, Algorithm 1, the cost table and the abstract would change.

\paragraph{Caveats}
\begin{itemize}
\item CIFAR-100-LT has two seeds, and the difference there is not resolved.
\item Only the two CIFAR benchmarks were run; the gate chose class rows on FER-2013, where the class means are close,
and there a two-moment sampler may do worse.
\item The sampler releases one more statistic: the $d\times d$ sum of $hh^\top$, $2560^2\approx6.6$ million numbers
per client on CIFAR. It is shared by all classes, so it reveals less per class than per-class covariances would, but
it is second-order information about every client's features. In these runs that sum was computed without secure
aggregation (\texttt{server.py}, \texttt{ssum(..., secure=False)}); a deployment would need it under secure
aggregation, at that size.
\item The ReLU after sampling makes the generated spread smaller than the exact one, as for the moment correction.
\end{itemize}

\end{document}
